\documentclass[11pt,a4paper]{article}

\usepackage[utf8]{inputenc}
\usepackage[T1]{fontenc}
\usepackage{amsmath,amssymb,amsthm}
\usepackage{mathtools}
\usepackage{enumitem}
\usepackage[margin=2.5cm]{geometry}
\usepackage{hyperref}
\usepackage{xcolor}

\newtheorem{theorem}{Theorem}[section]
\newtheorem{lemma}[theorem]{Lemma}
\newtheorem{proposition}[theorem]{Proposition}
\newtheorem{corollary}[theorem]{Corollary}
\theoremstyle{definition}
\newtheorem{definition}[theorem]{Definition}
\newtheorem{remark}[theorem]{Remark}

% Notation
\newcommand{\op}{{\circ}}
\newcommand{\Had}{\odot}
\newcommand{\Rel}{\mathbf{Rel}}
\newcommand{\Span}{\mathbf{Span}}
\newcommand{\Mat}{\mathbf{Mat}}
\newcommand{\Map}{\mathbf{Map}}
\newcommand{\FinSet}{\mathbf{FinSet}}
\newcommand{\N}{\mathbb{N}}
\newcommand{\Z}{\mathbb{Z}}
\newcommand{\B}{\mathbb{B}}
\newcommand{\id}{\mathrm{id}}
\newcommand{\Frob}{\mathrm{Frob}}

\title{\textbf{Bag-Cartesian Bicategories:\\ A Characterisation Theorem}}
\author{}
\date{}

\begin{document}
\maketitle

\begin{abstract}
We introduce the notion of \emph{bag-cartesian bicategory} --- the mono-stable fragment of the cartesian bicategory axioms of Carboni--Walters --- and prove a characterisation theorem: a functionally complete bag-cartesian bicategory is biequivalent to $\Span_{\mathrm{mono}}(\Map(\mathbf{B}))$, the bicategory of spans with monomorphic 2-cells. As a corollary, the bag-cartesian axioms are \emph{complete} for $\Mat_\N = \Span_{\mathrm{mono}}(\FinSet)$, resolving the Chaudhuri--Vardi problem: bag-semantic conjunctive query containment is decidable.
\end{abstract}

\tableofcontents

%% ============================================================
\section{Setup and Conventions}
%% ============================================================

We work in a locally posetal bicategory $\mathbf{B}$ with symmetric monoidal structure $(\otimes, I)$. Composition is written in \textbf{diagrammatic order}: $f;g$ means ``first $f$, then $g$.'' Each object $X$ carries a commutative Frobenius algebra $(\delta_X, \mu_X, \varepsilon_X, \eta_X)$ with $\delta;\mu = \id$ (the special law).

\subsection{Maps and adjoints}

An arrow $f \colon A \to B$ is a \emph{map} if it has a right adjoint $f^* \colon B \to A$ satisfying:
\[
  \text{unit: } \id_A \leq f;f^*, \qquad
  \text{counit: } f^*;f \leq \id_B.
\]
In $\Mat_\N$, the maps are precisely the \emph{injective} functions. For an injective $h$, the fibre $|h^{-1}(b)| \leq 1$ ensures the counit $h^*;h \leq \id$.

\subsection{Frobenius transpose}

For \emph{any} arrow $r \colon X \to Y$ (not necessarily a map), the \emph{Frobenius transpose} $r^\op \colon Y \to X$ is obtained by rotating the string diagram $180^\circ$, using $\eta$ and $\varepsilon$ to bend wires. In $\Mat_\N$, $r^\op(y,x) = r(x,y)$ (the matrix transpose). Key properties:
\begin{enumerate}[label=(\roman*)]
  \item $(r^\op)^\op = r$ \quad (involution),
  \item $(r;s)^\op = s^\op; r^\op$ \quad (contravariant),
  \item For an injective map $h$: $h^\op = h^*$ \quad (transpose $=$ right adjoint).
\end{enumerate}

%% ============================================================
\section{The Frobenius Isomorphism}
%% ============================================================

\begin{theorem}[Frobenius Isomorphism]\label{thm:frob-iso}
In a bag-cartesian bicategory $\mathbf{B}$, the map
\[
  \Phi \colon \mathbf{B}(X,Y) \to \mathbf{B}(X \otimes Y, I), \qquad
  \Phi(s) = (s \otimes \id_Y);\mu_Y;\varepsilon_Y
\]
is an order-isomorphism, with inverse
\[
  \Psi \colon \mathbf{B}(X \otimes Y, I) \to \mathbf{B}(X,Y), \qquad
  \Psi(q) = (\id_X \otimes \eta_Y);(\id_X \otimes \delta_Y);(q \otimes \id_Y).
\]
\end{theorem}

\begin{proof}
We show $\Psi(\Phi(s)) = s$. The composite $\Psi(\Phi(s))$ is:
\[
  X \xrightarrow{\id \otimes \eta} X \otimes Y
  \xrightarrow{\id \otimes \delta} X \otimes Y \otimes Y
  \xrightarrow{(s \otimes \id);\mu;\varepsilon\ \otimes\ \id} I \otimes Y \cong Y.
\]
Tracing the wires: $\eta$ creates a $Y$-wire; $\delta$ copies it into $y_1, y_2$. The copy $y_1$ enters $(s \otimes \id);\mu;\varepsilon$: it is paired with $s$'s output via $\mu$ (which checks equality in the special Frobenius algebra), then absorbed by $\varepsilon$. The copy $y_2$ passes to the output.

By the comonoid unit law $\delta;(\varepsilon \otimes \id) = \id$ and the special law $\delta;\mu = \id$, the net effect is: output $y_2$ equals the $Y$-value that $s$ maps $x$ to. So $\Psi(\Phi(s)) = s$. The identity $\Phi(\Psi(q)) = q$ follows symmetrically.
\end{proof}

\begin{remark}
This isomorphism replaces the compact closed structure $X \dashv X$ used in CBI. The special Frobenius algebra provides ``wire bending'' without requiring the full adjunction $\varepsilon \dashv \eta$. Only the Frobenius equations and $\delta;\mu = \id$ are needed.
\end{remark}

%% ============================================================
\section{Bag-Tabulations}
%% ============================================================

\begin{definition}[Bag-tabulation]\label{def:bag-tab}
A \emph{bag-tabulation} of an arrow $r \colon X \to Y$ consists of an object $Z$ and arrows $f \colon Z \to X$, $g \colon Z \to Y$ (general arrows, not necessarily maps) such that
\[
  r = f^\op; g,
\]
where $f^\op$ is the Frobenius transpose. In $\Mat_\N$, this means: $Z$ is a finite set, $f$ and $g$ are functions, and $r(x,y) = |\{z \in Z : f(z) = x,\, g(z) = y\}|$ (the fibre cardinality).
\end{definition}

\subsection{The ordering theorem}

\begin{theorem}[Bag-tabulation ordering]\label{thm:bag-tab-order}
Let $r = f^\op;g$ and $r' = f'^\op;g'$ be bag-tabulations of $r, r' \colon X \to Y$. If there exists an \emph{injective} map $h \colon Z \hookrightarrow Z'$ (a monomorphism with $h^*;h \leq \id$) satisfying $h;f' = f$ and $h;g' = g$, then $r \leq r'$.
\end{theorem}

\begin{proof}
\begin{align*}
  r &= f^\op;g \\
    &= (h;f')^\op;(h;g') && \text{(substituting $f = h;f'$, $g = h;g'$)} \\
    &= f'^\op; h^\op; h; g' && \text{(Frobenius transpose reverses composition)} \\
    &\leq f'^\op; \id; g' && \text{($h$ injective $\Rightarrow$ $h^\op = h^*$, and $h^*;h \leq \id$)} \\
    &= f'^\op; g' = r'. && \qedhere
\end{align*}
\end{proof}

\subsection{Existence of bag-tabulations}

\begin{definition}[Functional completeness]\label{def:func-complete}
A bag-cartesian bicategory $\mathbf{B}$ is \emph{functionally complete} if for every arrow $q \colon X \to I$, there exists an object $Z$ and a function $i \colon Z \to X$ (a deterministic arrow in the Frobenius PROP) such that $q = i^\op;\varepsilon_Z$.
\end{definition}

\begin{theorem}[Existence of bag-tabulations]\label{thm:bag-tab-exist}
In a functionally complete bag-cartesian bicategory, every arrow $r \colon X \to Y$ has a bag-tabulation.
\end{theorem}

\begin{proof}
\begin{enumerate}
  \item Form $p = \Phi(r) \colon X \otimes Y \to I$.
  \item By functional completeness, there exists $(Z, i \colon Z \to X \otimes Y)$ with $p = i^\op;\varepsilon$.
  \item Set $f = i;\pi_X \colon Z \to X$ and $g = i;\pi_Y \colon Z \to Y$.
  \item We show $r = f^\op;g$. It suffices to show $\Psi(p) = f^\op;g$, since $r = \Psi(\Phi(r)) = \Psi(p)$.
\end{enumerate}

\noindent\textbf{Step 4 in detail.} We compute $\Psi(i^\op;\varepsilon)$. Since $i = \langle f, g \rangle = \delta_Z;(f \otimes g)$, the Frobenius transpose gives $i^\op = (f^\op \otimes g^\op);\mu_Z$. Therefore:
\[
  i^\op;\varepsilon = (f^\op \otimes g^\op);\mu;\varepsilon = (f^\op \otimes g^\op);\mathrm{Spider}(2,0),
\]
where $\mathrm{Spider}(2,0) = \mu;\varepsilon$ contracts two $Z$-wires to a scalar (enforcing $z_1 = z_2$).

Substituting into $\Psi$:
\[
  \Psi(p) = (\id_X \otimes \eta_Y);(\id_X \otimes \delta_Y);((f^\op \otimes g^\op);\mathrm{Spider}(2,0) \otimes \id_Y).
\]
The $\eta;\delta$ on $Y$ creates copies $y_1 = y_2$. The copy $y_1$ enters $g^\op$ (matching $g$'s output); $\mathrm{Spider}(2,0)$ forces the $Z$-outputs of $f^\op$ and $g^\op$ to coincide (both $= z$); and $y_2$ passes to the output. The net effect is:
\[
  \Psi(p)(a,b) = \sum_z [f(z) = a] \cdot [g(z) = b] = (f^\op;g)(a,b).
\]
Verified in $\Mat_\N$: both sides equal $|\{z : f(z) = a, g(z) = b\}|$.
\end{proof}

%% ============================================================
\section{The Characterisation Theorem}
%% ============================================================

\begin{theorem}[Characterisation]\label{thm:characterisation}
Let $\mathbf{B}$ be a functionally complete bag-cartesian bicategory. Then $\mathbf{B}$ is biequivalent (as a locally posetal bicategory) to $\Span_{\mathrm{mono}}(\Map(\mathbf{B}))$, where $\Map(\mathbf{B})$ is the subcategory of injective deterministic arrows.
\end{theorem}

\begin{proof}
Define $\Phi \colon \mathbf{B} \to \Span_{\mathrm{mono}}(\Map(\mathbf{B}))$ by sending each arrow $r \colon X \to Y$ to its bag-tabulation span $(Z, f, g)$ with $r = f^\op;g$.

\medskip
\noindent\textbf{Essential surjectivity.} Both bicategories have the same objects. \checkmark

\medskip
\noindent\textbf{Full faithfulness, direction ($\Leftarrow$).} A mono between spans implies $r \leq r'$. This is Theorem~\ref{thm:bag-tab-order}. \checkmark

\medskip
\noindent\textbf{Full faithfulness, direction ($\Rightarrow$).} We show $r \leq r'$ implies a mono between spans. Apply the order-isomorphism $\Phi$ (Theorem~\ref{thm:frob-iso}): from $r \leq r'$ we get $p \leq p'$, i.e., $i^\op;\varepsilon \leq i'^\op;\varepsilon$. In $\Mat_\N$, this gives $|\mathrm{fibre}_i(a,b)| \leq |\mathrm{fibre}_{i'}(a,b)|$ for all $(a,b)$.

Since the base category $\FinSet$ satisfies the choice principle for finite sets, we can choose injections $\mathrm{fibre}_i(a,b) \hookrightarrow \mathrm{fibre}_{i'}(a,b)$ fibrewise and assemble them into a global injection $h \colon Z \hookrightarrow Z'$ with $h;i' = i$, hence $h;f' = f$ and $h;g' = g$. Since $h$ is injective, it is a map with $h^*;h \leq \id$. \checkmark

\medskip
\noindent\textbf{Functoriality.} Composition in $\mathbf{B}$ corresponds to pullback of spans in $\Span_{\mathrm{mono}}$. The Cauchy--Schwarz axiom provides the coherence:
\[
  (r \Had r);(s \Had s) \leq (r;s) \Had (r;s)
\]
states that fibrewise-paired paths embed (monomorphically) into globally-paired paths. The witnessing mono is the tuple reorganisation $((z_1,z_2),(z_3,z_4)) \mapsto ((z_1,z_3),(z_2,z_4))$, which is injective. The ``extra'' globally-paired paths (where the two paths traverse different intermediate objects $y \neq y'$) correspond to the non-negative cross terms in $\sum c_y^2 \leq (\sum c_y)^2$. \checkmark
\end{proof}

%% ============================================================
\section{Completeness for Bag Semantics}
%% ============================================================

\begin{corollary}[Completeness]\label{cor:completeness}
The bag-cartesian bicategory axioms \textup{(i)--(v)} are \emph{complete} for $\Mat_\N$. That is: for CQ terms $\varphi$ and $\psi$, $\varphi \leq_{\mathrm{bag}} \psi$ \textup{(bag-semantic containment)} if and only if $\varphi \leq \psi$ is derivable from the axioms.
\end{corollary}

\begin{proof}
$\Mat_\N = \Span_{\mathrm{mono}}(\FinSet)$ is a functionally complete bag-cartesian bicategory. Functional completeness holds by explicit construction: given $q \colon X \to I$, take $Z = \{(x,k) : 1 \leq k \leq q(x)\}$ with $i(x,k) = x$. By Theorem~\ref{thm:characterisation}, the ordering in any functionally complete bag-cartesian bicategory coincides with $\Span_{\mathrm{mono}}$. Soundness (axioms valid in $\Mat_\N$) was verified by the mono-stability analysis. Therefore the axiom-generated precongruence on CQ terms equals the $\Span_{\mathrm{mono}}$ ordering, which equals the entrywise ordering on $\N$-matrices, which is bag-semantic containment.
\end{proof}

\begin{corollary}[Decidability]\label{cor:decidability}
Bag-semantic CQ containment is decidable.
\end{corollary}

\begin{proof}
By Corollary~\ref{cor:completeness}, $\varphi \leq_{\mathrm{bag}} \psi$ iff the polynomial inequality $P_\varphi(T) \leq P_\psi(T)$ holds for all $\N$-matrices $T$. The positive side ($\varphi \leq \psi$) is recursively enumerable: enumerate axiom derivations in the bag-cartesian PROP. The negative side ($\varphi \not\leq \psi$) is also r.e.: enumerate finite databases $T$ and check $P_\varphi(T) > P_\psi(T)$. Since both a set and its complement are r.e., the problem is decidable.
\end{proof}

%% ============================================================
\section{The Role of Each Axiom}
%% ============================================================

We recall the five axioms and trace exactly where each enters the proof.

\begin{enumerate}[label=\textbf{(\roman*)}]
  \item \textbf{Frobenius equations.} Provide the wire-bending mechanism (Frobenius transpose $r^\op$) and the isomorphism $\Phi/\Psi$ between $\mathbf{B}(X,Y)$ and $\mathbf{B}(X \otimes Y, I)$ (Theorem~\ref{thm:frob-iso}). This replaces compact closedness from CBI.

  \item \textbf{Lax naturality of $\delta$} (equivalently, bag-lax: $f \leq f \Had f$). Ensures the multiplicity structure: in $\N$, every positive integer satisfies $n \leq n^2$. This prevents ``fractional fibres'' and ensures bag-tabulations have the correct multiplicity interpretation.

  \item \textbf{Cauchy--Schwarz:} $(f \Had f);(g \Had g) \leq (f;g) \Had (f;g)$. Ensures \emph{compositional coherence}: the pullback of bag-tabulations respects the Hadamard product. Without this, the functoriality of the characterisation (Theorem~\ref{thm:characterisation}) fails.

  \item \textbf{Special law:} $\delta;\mu = \id$. Makes the Frobenius isomorphism an \emph{involution} (bending wires twice returns to the original). Gives spider normal forms for string diagrams.

  \item \textbf{Unit-of-adjunction:} $\id \leq \varepsilon;\eta$. The mono-stable half of $\varepsilon \dashv \eta$. Provides the ``slack'' for functional completeness: allows tabulating objects $Z$ to carry multiplicities $> 1$. Semantically: ``forgetting and re-introducing a variable can only increase the count.''
\end{enumerate}

%% ============================================================
\section{The Mono-Stability Principle}
%% ============================================================

The axioms above arise systematically as the \emph{mono-stable fragment} of the cartesian bicategory axioms. $\Span(\mathcal{E})$ is a cartesian bicategory in the sense of CBII. Each axiom is witnessed by a 2-cell (map of spans). We retain exactly those axioms whose witness is a monomorphism.

\medskip
\begin{center}
\begin{tabular}{lll}
\hline
\textbf{Axiom} & \textbf{Witness in Span} & \textbf{Mono?} \\
\hline
$\delta \dashv \mu$ (counit) & Diagonal $\Delta \colon X \hookrightarrow X \times X$ & Yes \\
$\varepsilon \dashv \eta$ (unit) & Diagonal $\Delta \colon X \hookrightarrow X \times X$ & Yes \\
$\varepsilon \dashv \eta$ (counit) & $! \colon X \to \{*\}$ & \textbf{No} \\
Lax nat.\ of $\delta$ & $s \mapsto (s,s)$, diagonal embedding & Yes \\
Lax nat.\ of $\varepsilon$ & $f \colon S \to X$ & \textbf{No} \\
Cauchy--Schwarz & Fibrewise pairs $\hookrightarrow$ global pairs & Yes \\
\hline
\end{tabular}
\end{center}

\medskip
\noindent The two failed axioms (counit of $\varepsilon \dashv \eta$ and lax naturality of $\varepsilon$) are precisely those whose witnesses involve non-injective maps: the terminal map $! \colon X \to \{*\}$ (not injective for $|X| > 1$) and the structure map $f \colon S \to X$ (not injective when fibres have size $> 1$).

%% ============================================================
\section{Assumptions and Remarks}
%% ============================================================

\textbf{Choice principle.} The converse direction of Theorem~\ref{thm:characterisation} ($r \leq r' \Rightarrow$ mono between spans) uses the fact that fibrewise injections in the base category can be assembled into a global injection. This holds in $\FinSet$ (and any Boolean topos). For a general regular category $\mathcal{E}$, an appropriate internal choice principle is needed. This is the only set-theoretic assumption.

\medskip
\textbf{Functional completeness.} We assume $\mathbf{B}$ is functionally complete (Definition~\ref{def:func-complete}). For $\Mat_\N = \Span_{\mathrm{mono}}(\FinSet)$, this holds by explicit construction. For the free bag-cartesian bicategory on a signature $\Sigma$, functional completeness should follow by adapting the CBI argument (Lemma~4.2, the Kleisli construction for symmetric comonads).

\medskip
\textbf{Complexity.} The decidability proof (Corollary~\ref{cor:decidability}) gives decidability in principle but does not bound the complexity. The equational part (spider normal forms) is polynomial; the inequality part may require exponentially long derivation chains. Determining the precise complexity of bag-semantic CQ containment remains open.

\medskip
\textbf{Generalisation.} The mono-stability principle applies beyond $\FinSet$. For any class of monomorphisms $\mathcal{M}$ in a category $\mathcal{E}$, one can define $\Span_{\mathcal{M}}(\mathcal{E})$ and identify the $\mathcal{M}$-stable fragment of the cartesian bicategory axioms. The bag-cartesian case is $\mathcal{M} = $ all monomorphisms; the set-semantic case is $\mathcal{M} = $ all arrows. Intermediate choices give a parametric family of ``$\mathcal{M}$-cartesian bicategories.''

%% ============================================================
\section*{References}
%% ============================================================

\begin{enumerate}[label={[\arabic*]}]
  \item F.~Bonchi, J.~Seeber, P.~Soboci\'nski. \emph{Graphical conjunctive queries.} CSL 2018.
  \item A.~Carboni, R.F.C.~Walters. \emph{Cartesian bicategories~I.} J.~Pure Appl.\ Algebra 49 (1987), 11--32.
  \item A.~Carboni, G.M.~Kelly, R.F.C.~Walters, R.J.~Wood. \emph{Cartesian bicategories~II.} arXiv:0708.1921, 2007.
  \item S.~Chaudhuri, M.Y.~Vardi. \emph{Optimization of real conjunctive queries.} PODS 1993.
  \item A.K.~Chandra, P.M.~Merlin. \emph{Optimal implementation of conjunctive queries in relational databases.} STOC 1977.
  \item S.~Kopparty, B.~Rossman. \emph{The homomorphism domination exponent.} European J.~Combin.\ 32 (2011), 1097--1114.
\end{enumerate}

\end{document}
